\documentclass[12pt]{article}
\usepackage[utf8]{inputenc}
\usepackage[T1]{fontenc}
\usepackage[margin=1in]{geometry}
\usepackage{setspace}
\usepackage{xcolor}
\usepackage[colorlinks=true, urlcolor=blue!60!black, linkcolor=blue!60!black]{hyperref}
\setlength{\parindent}{2em}
\setlength{\parskip}{1.25em}
\renewcommand{\baselinestretch}{1.0}
\title{\textit{INS v. Chadha}, 462 U.S. 919 (1983)}
\author{Case Brief}
\date{}

\begin{document}
\maketitle

\section{Case Metadata}
\begin{description}
  \item[Citation:] \textit{Immigration and Naturalization Service v. Chadha}, 462 U.S. 919 (1983) (No.~80-1832; consolidated with No.~80-2170, \textit{United States House of Representatives v. INS}, and No.~80-2171, \textit{United States Senate v. INS}).
  \item[Argued / Decided:] Argued February 22, 1982; Reargued December 7, 1982; Decided June 23, 1983.
  \item[Vote:] 7--2 (affirming the U.S. Court of Appeals for the Ninth Circuit, 634 F.2d 408 (9th Cir. 1980)).
  \item[Author:] Chief Justice Warren E. Burger (joined by Justices Brennan, Marshall, Blackmun, Stevens, and O'Connor).
  \item[Concurrences / Dissents:] Justice Powell (concurring in the judgment, 462 U.S. at 959); Justice White (dissenting, 462 U.S. at 967); Justice Rehnquist (dissenting, joined by Justice White, 462 U.S. at 1013).
  \item[Chapter / Topic:] Separation of Powers and the Executive --- Bicameralism, Presentment, and the Legislative Veto (U.S. Const. art.~I, \S\S~1, 7).
  \item[Primary Source URLs:] \href{https://supreme.justia.com/cases/federal/us/462/919/}{Justia U.S. Supreme Court Center} (\url{https://supreme.justia.com/cases/federal/us/462/919/}); \href{https://www.law.cornell.edu/supremecourt/text/462/919}{Cornell LII} (\url{https://www.law.cornell.edu/supremecourt/text/462/919}).
\end{description}

\section{Facts and Procedural History}
Jagdish Rai Chadha, an East Indian born in Kenya holding a British passport, was lawfully admitted to the United States in 1966 on a nonimmigrant student visa. After his visa expired on June 30, 1972, the Immigration and Naturalization Service (INS) ordered him in October 1973 to show cause why he should not be deported. At his January 1974 deportation hearing, Chadha conceded deportability but applied for suspension of deportation under \S~244(a)(1) of the Immigration and Nationality Act (INA), 8 U.S.C. \S~1254(a)(1), which authorized the Attorney General to suspend deportation for an alien continuously present in the United States for at least seven years, of good moral character, whose deportation would result in ``extreme hardship.''

On June 25, 1974, an Immigration Judge found that Chadha satisfied all statutory criteria, ordered his deportation suspended, and transmitted a report to Congress as required by \S~244(c)(1). Under \S~244(c)(2) of the INA, 8 U.S.C. \S~1254(c)(2), if either the House of Representatives or the Senate passed a resolution disapproving the suspension during the session in which it was reported or the next session, the Attorney General was required to deport the alien. On December 16, 1975---three days before the statutory window expired---the House passed H.R. Res. 926 without printed text, debate, or a recorded vote, opposing permanent residence for Chadha and five other aliens out of 340 reported suspensions. The resolution was neither sent to the Senate nor presented to the President.

The Immigration Judge reopened proceedings and ordered Chadha deported on November 8, 1976, holding he lacked authority to pass on the constitutionality of \S~244(c)(2); the Board of Immigration Appeals dismissed Chadha's appeal on the same ground. Chadha petitioned for review in the Ninth Circuit under \S~106(a) of the INA. The INS agreed with Chadha that the one-House veto was unconstitutional, prompting the Ninth Circuit to invite the House and Senate to participate as amici (and later as formal intervenors). The Ninth Circuit held \S~244(c)(2) unconstitutional under the separation of powers and directed the Attorney General to cease deportation efforts. The INS appealed under 28 U.S.C. \S~1252, and both Houses of Congress petitioned for certiorari.

\section{Issue(s) Presented}
\begin{enumerate}
  \item \textbf{Threshold Justiciability and Severability:} Did the Supreme Court have appellate jurisdiction under 28 U.S.C. \S~1252 when the INS agreed with Chadha on the merits; was \S~244(c)(2) severable from the remainder of \S~244; did Chadha have Article III standing despite speculative alternative avenues of relief; and did the challenge present a nonjusticiable political question?
  \item \textbf{Constitutional Merits:} Does a statutory provision (\S~244(c)(2) of the INA) authorizing one House of Congress, by simple resolution without passage by the other House or presentment to the President, to invalidate an Executive Branch decision made pursuant to congressionally delegated authority violate the Bicameralism and Presentment Clauses of Article I, \S\S~1 and 7?
\end{enumerate}

\section{Holding}
Yes on justiciability and severability, and yes on the merits. The Court held that \S~244(c)(2) is severable from the remainder of \S~244, that the dispute presented a justiciable Article III case or controversy, and that the one-House legislative veto in \S~244(c)(2) is unconstitutional because it constitutes an exercise of legislative power that fails to comply with the bicameralism and presentment requirements of Article I, \S\S~1 and 7.

\section{Reasoning}
Chief Justice Burger first cleared seven threshold objections (Part II): (1) the INS was an ``aggrieved'' party under 28 U.S.C. \S~1252 because the Ninth Circuit enjoined it from carrying out a deportation order it otherwise would have executed; (2) \S~244(c)(2) was presumed severable under the INA's express severability clause (\S~406), legislative history showing Congress adopted \S~244 to escape the intolerable burden of private immigration bills, and the fact that \S~244 remained a ``fully operative'' report-and-wait mechanism without the veto; (3) Chadha had Article III standing because invalidating the veto redressed his threatened deportation; (4) speculative alternative relief (such as Chadha's 1980 marriage to a U.S. citizen or the Refugee Act of 1980) did not bar adjudication; (5) the Ninth Circuit had jurisdiction under \S~106(a) because the final deportation order was contingent on the House veto; (6) concrete adverseness existed both before and after Congress intervened; and (7) no political question under \textit{Baker v. Carr}, 369 U.S. 186 (1962), barred judicial review of whether Congress chose constitutionally permissible means to exercise its Naturalization Clause power.

On the merits (Parts III--IV), Burger adopted a rigorous formalist reading of Article I:
\begin{itemize}
  \item \textbf{Efficiency Does Not Override Constitutional Structure:} Although legislative vetoes had proliferated (295 provisions in 196 statutes since 1932), convenience and utility cannot save a procedure that contravenes explicit constitutional text.
  \item \textbf{Bicameralism and Presentment as an Integrated Design:} Article I, \S~1 (vesting all legislative powers in a bicameral Congress) and Article I, \S~7, cls. 2--3 (requiring every bill, order, resolution, or vote requiring concurrence of both Houses to be presented to the President, subject to veto and two-thirds override) establish a ``single, finely wrought and exhaustively considered, procedure.'' Presentment protects the Executive from legislative encroachment, guards against improvident laws, and injects a national perspective; bicameralism divides legislative power to prevent legislative despotism and protect both large and small states.
  \item \textbf{Legislative Purpose and Effect Test:} Not every congressional act requires bicameralism and presentment, but an action is legislative if it has ``the purpose and effect of altering the legal rights, duties, and relations of persons, including the Attorney General, Executive Branch officials and Chadha, all outside the Legislative Branch.'' Here, absent \S~244(c)(2), Congress could compel Chadha's deportation only by enacting a statute; ``Congress must abide by its delegation of authority until that delegation is legislatively altered or revoked.''
  \item \textbf{Exhaustive Unicameral Exceptions:} Where the Framers intended one House to act alone with binding force, they specified those exceptions narrowly and explicitly: House impeachment (Art.~I, \S~2, cl.~5), Senate impeachment trials (Art.~I, \S~3, cl.~6), Senate confirmation of appointments (Art.~II, \S~2, cl.~2), and Senate treaty ratification (Art.~II, \S~2, cl.~2). Expressio unius forecloses implied unicameral vetoes.
  \item \textbf{Distinguishing Executive Rulemaking (Footnote 16):} Responding to Congress's claim that the Attorney General's suspension was itself unpresented ``lawmaking,'' Burger reasoned that executive action under a duly enacted statute is presumptively Article II execution, bounded by statutory standards, judicial review, and Congress's power to amend or revoke the delegation.
\end{itemize}

\section{Concurring and Dissenting Opinions}
\begin{itemize}
  \item \textbf{Justice Powell (Concurring in the Judgment, 462 U.S. at 959--967):} Warned that the majority's sweeping Presentment Clause holding would invalidate nearly every legislative veto across hundreds of statutes. Powell urged deciding the case on a narrower separation-of-powers ground: by determining that six named individuals failed to satisfy statutory hardship criteria---without general rules, counsel, or a hearing---the House assumed an adjudicatory, judicial function analogous to a bill of attainder, exposing individual rights to ``the tyranny of a shifting majority.''
  \item \textbf{Justice White (Dissenting, 462 U.S. at 967--1013):} Delivered a landmark functionalist dissent (with a 56-statute appendix). White argued that the legislative veto was an ``indispensable, political invention'' allowing Congress to delegate necessary authority to the modern administrative state without abdicating democratic accountability, and that invalidating it forced Congress into a ``Hobson's choice'' between impossible statutory specificity and wholesale surrender to unelected agencies. Substantively, White contended that: (1) agencies routinely make binding ``quasi-legislative'' rules without bicameralism and presentment, so Congress should be permitted to reserve a check on its own delegation; (2) under \S~244, the Attorney General's suspension was merely a legislative proposal, and a departure from the status quo (deportability) still required the concurrence of the President (via the Attorney General) and both Houses (by silence); and (3) under \textit{Nixon v. Administrator of General Services}, 433 U.S. 425 (1977), \S~244(c)(2) did not prevent the Executive from accomplishing any constitutionally assigned function.
  \item \textbf{Justice Rehnquist (Dissenting, joined by Justice White, 462 U.S. at 1013--1016):} Focused solely on severability. Reviewing the legislative history of the INA, Rehnquist argued that Congress never intended to grant the Attorney General unreviewable discretion to cancel deportations; because the one-House veto was inseverable from \S~244(a)(1), striking the veto destroyed the Attorney General's suspension power as well, depriving Chadha of redressable relief and standing.
\end{itemize}

\section{Key Quotations}
\begin{quote}
``Convenience and efficiency are not the primary objectives---or the hallmarks---of democratic government\ldots'' \\
--- Chief Justice Burger, 462 U.S. at 944.
\end{quote}

\begin{quote}
``It emerges clearly that the prescription for legislative action in Art. I, \S\S~1, 7, represents the Framers' decision that the legislative power of the Federal Government be exercised in accord with a single, finely wrought and exhaustively considered, procedure.'' \\
--- Chief Justice Burger, 462 U.S. at 951.
\end{quote}

\begin{quote}
``In purporting to exercise power defined in Art. I, \S~8, cl. 4, to `establish an uniform Rule of Naturalization,' the House took action that had the purpose and effect of altering the legal rights, duties, and relations of persons, including the Attorney General, Executive Branch officials and Chadha, all outside the Legislative Branch.'' \\
--- Chief Justice Burger, 462 U.S. at 952.
\end{quote}

\begin{quote}
``With all the obvious flaws of delay, untidiness, and potential for abuse, we have not yet found a better way to preserve freedom than by making the exercise of power subject to the carefully crafted restraints spelled out in the Constitution.'' \\
--- Chief Justice Burger, 462 U.S. at 959.
\end{quote}

\begin{quote}
``Today's decision strikes down in one fell swoop provisions in more laws enacted by Congress than the Court has cumulatively invalidated in its history.'' \\
--- Justice White, dissenting, 462 U.S. at 1002.
\end{quote}

\section{Doctrinal Significance and Modern Legacy (Through 2026)}
\textit{Chadha} is the foundational pillar of modern separation-of-powers formalism. Two weeks later, the Court summarily affirmed D.C. Circuit rulings invalidating a two-House veto over independent agency rulemaking in \textit{Process Gas Consumers Group v. Consumer Energy Council of America}, 463 U.S. 1216 (1983), confirming that \textit{Chadha} bars one-House, two-House, and committee vetoes alike. The Court extended \textit{Chadha}'s anti-circumvention logic in \textit{Bowsher v. Synar}, 478 U.S. 714 (1986) (barring Congress from assigning executive budget cuts to a legislatively removable Comptroller General), \textit{Alaska Airlines, Inc. v. Brock}, 480 U.S. 678 (1987) (holding legislative vetoes presumptively severable), and \textit{Metropolitan Washington Airports Authority v. Citizens for Abatement of Aircraft Noise, Inc.}, 501 U.S. 252 (1991) (invalidating a board of review staffed by Members of Congress acting in ``individual capacities'').

Through 2026, \textit{Chadha}'s structural consequences continue to shape executive-legislative relations in four major ways, as cataloged in the Congressional Research Service's June 2026 report, \href{https://www.everycrsreport.com/reports/R48965.html}{\textit{Statutory Mechanisms for Agency Oversight After INS v. Chadha} (CRS Report R48965)}:
\begin{enumerate}
  \item \textbf{Statutory Workarounds and the Congressional Review Act (CRA):} Unable to use concurrent resolutions, Congress turned to report-and-wait mechanisms paired with fast-track joint resolutions of disapproval requiring presidential signature---most prominently the Congressional Review Act of 1996, 5 U.S.C. \S\S~801--808. Because a sitting President will almost always veto a disapproval of his own administration's rules, the CRA functions effectively only during presidential transitions under its ``lookback'' mechanism, prompting renewed 2025--2026 scholarly and legislative debates over reforming the CRA (see, e.g., \href{https://harvardlawreview.org/blog/2026/05/from-destruction-to-construction-the-case-for-a-new-congressional-review-act/}{Harvard Law Review Blog, May 2026}).
  \item \textbf{The REINS Act Debate:} Proposals such as the Regulations from the Executive in Need of Scrutiny (REINS) Act of 2025 (H.R. 142, 119th Cong.) would require affirmative bicameral approval and presentment before any ``major rule'' takes effect. Defenders argue it satisfies \textit{Chadha} by treating agency drafts as proposals; critics contend that allowing one House's inaction to block a rule functions as a one-House veto in disguise.
  \item \textbf{Informal Persistence of Committee Vetoes:} Despite \textit{Chadha}, Congress has enacted hundreds of post-1983 provisions requiring committee approval (including in the Consolidated Appropriations Act, 2026, Pub. L. No. 119-75), relying on informal appropriations leverage rather than judicial enforceability.
  \item \textbf{Asymmetry with Broad Delegations:} \textit{Chadha} eliminated Congress's post-hoc veto while leaving the permissive ``intelligible principle'' nondelegation doctrine intact---a framework the Supreme Court reaffirmed 6--3 in \href{https://ballotpedia.org/Federal_Communications_Commission_v._Consumers%E2%80%99_Research}{\textit{FCC v. Consumers' Research} (June 27, 2025)}. Critics argue this combination---amplified in emergency-powers statutes like the National Emergencies Act, where terminating a presidential emergency now requires a veto-proof two-thirds supermajority---has produced a ``\textit{Chadha} Presidency'' tilting structural power toward the Executive.
\end{enumerate}

\section{Pedagogical Critiques and Exam Hypotheticals}
\textbf{Pedagogical Critiques:}
\begin{itemize}
  \item \textit{The Footnote 16 Asymmetry:} Students should interrogate Burger's claim that agency rulemaking is merely ``executive'' while a congressional veto of that rule is ``legislative.'' If both alter ``the legal rights, duties, and relations of persons\ldots outside the Legislative Branch,'' \textit{Chadha}'s test struggles to explain why the Executive can make binding rules without bicameralism and presentment, whereas Congress cannot conditionally reserve a check on them.
  \item \textit{Severability Distortion:} By severing \S~244(c)(2) rather than striking \S~244 altogether, the Court arguably gave the Executive broader, unchecked power than Congress ever agreed to delegate---the exact concern driving Rehnquist's dissent.
\end{itemize}

\textbf{Exam Hypothetical:} Congress enacts the Artificial Intelligence Safety Act of 2026, authorizing the Secretary of Commerce to issue binding national safety standards for frontier AI models. Section 12 provides that any standard shall take effect 60 days after submission to Congress unless both the House and Senate pass a concurrent resolution of disapproval within that window. Section 13 adds that if \S~12 is held invalid, the Secretary's rulemaking authority under the Act shall automatically terminate. The Secretary issues a binding rule; both Houses pass a concurrent resolution disapproving it without presenting it to the President. An AI developer sues to enforce the disapproval (or a safety group sues to enforce the rule). \textit{Analysis:} Under \textit{Chadha} and \textit{Process Gas}, even a \emph{two-House} concurrent resolution violates Presentment (Art.~I, \S~7, cl.~3) because it alters legal rights and duties outside Congress without presidential presentment. However, unlike \textit{Chadha} and \textit{Alaska Airlines}, \S~13 contains an express \emph{inseverability} (fallback) clause, meaning a court must honor Congress's explicit text and invalidate the Secretary's underlying rulemaking authority as well.

\textbf{``Opposite Side'' Contrast:} Compare \textit{Chadha} with \textit{Clinton v. City of New York}, 524 U.S. 417 (1998), and \textit{Mistretta v. United States}, 488 U.S. 361 (1989). Whereas \textit{Chadha} applies strict Article I, \S~7 formalism to block \emph{Congress} from bypassing bicameralism and presentment to claw back delegated authority, \textit{Clinton} applies the exact same ``finely wrought'' formalism in reverse to block the \emph{President} from unilaterally canceling enacted statutory provisions---even with Congress's statutory permission. By contrast, \textit{Mistretta} (and \textit{FCC v. Consumers' Research} (2025)) illustrates the Court's functionalist tolerance when Congress delegates broad rulemaking authority outward to an agency or commission rather than retaining a post-enactment legislative control mechanism.

\end{document}
